\section{Introduction}

Not all execution feedback is created equal: fine-grained rewards that pinpoint the exact error location can backfire when that localization is unreliable. This counterintuitive observation motivates our work on error-type-gated feedback for reinforcement learning (RL) fine-tuning of code language models.

\subsection{Background}

Recent advances in RL fine-tuning have enabled code LLMs to improve through execution feedback. Methods like CodeRL~\cite{le2022coderl} and RLTF~\cite{liu2023rltf} leverage program execution outcomes to provide reward signals, with RLTF demonstrating that combining coarse (pass/fail) and fine-grained (token-level) feedback outperforms single-signal approaches. VeRPO~\cite{wang2026verpo} further showed that how feedback is aggregated matters---cardinality bias in dense rewards can degrade performance. These findings establish that feedback granularity and aggregation strategy are critical design choices.

\subsection{The Problem}

However, existing methods apply fine-grained feedback uniformly, regardless of whether the error localization is reliable. When a Python program raises a SyntaxError, the traceback points to the exact buggy token. But when it raises a RuntimeError---say, due to a mishandled edge case---the traceback often points to a symptom line far from the root cause. Applying token-level reward penalties at such misleading locations injects noise into gradient updates: the model receives credit assignment signals at wrong tokens, potentially learning incorrect corrections.

\subsection{Key Insight}

We observe that Python exception types naturally partition into two categories with dramatically different localization reliability. U\_line errors (SyntaxError, NameError, TypeError, etc.) have 100\% traceback accuracy---the reported line matches the actual bug location. U\_ignore errors (RuntimeError, RecursionError, MemoryError, etc.) have only 20\% accuracy---tracebacks point to symptom locations. This 80-percentage-point gap, which we measure empirically (chi-square $p=9.57\times10^{-74}$), suggests that conditioning feedback strategy on error type could substantially reduce gradient noise.

\subsection{Our Approach}

We propose error-type-gated fine-grained feedback: apply token-level penalties only for errors where localization is reliable (U\_line category), while falling back to coarse-only feedback for unreliably-localized errors (U\_ignore category). This simple modification preserves the benefits of fine-grained credit assignment where it is accurate while avoiding noise injection where it would be misleading.

\subsection{Contributions}

Through a systematic verification pipeline, we establish the causal mechanism underlying this approach. We first confirm that fine-grained penalties do concentrate gradients at traceback locations (16.11$\times$ concentration ratio). We then demonstrate that this localization varies dramatically by error type (100\% vs 20\% accuracy). We show that unreliable localization causes measurable gradient noise at ground-truth locations ($p<10^{-13}$). Finally, we demonstrate that gating improves signal-to-noise ratio (+4.61\%) and training efficiency (gated condition reaches 30\% pass@1 while baseline does not).

Our work introduces credit assignment reliability as a lens for understanding feedback granularity in code RL, providing both theoretical framing and empirical validation for error-type-aware training strategies.
